\section{Introduction}
\label{sec:introduction}

Configurable software systems---compilers, video encoders, databases, garbage
collectors, archivers---expose tens to hundreds of configuration options whose
joint effect on non-functional properties is rarely
linear~\cite{kaltenecker2019,julianasurvey,siegmund2015FSE}. Predicting the
performance of an arbitrary configuration without measuring it is therefore a
long-standing problem, typically approached by \emph{sampling, measuring,
learning}: a sampling strategy selects a subset of configurations, those
configurations are executed and measured, a regression model is trained on the
resulting (configuration, performance) pairs, and the model is then evaluated
on the rest of the configuration space~\cite{guo2013ASE,sarkar2015,kaltenecker2019,siegmund2015FSE,nair2018c}.
Because measuring all configurations is usually
infeasible~\cite{siegmund2015FSE,julianasurvey}, sampling strategies are
commonly compared along two axes: prediction accuracy and \emph{measurement
effort}.

\paragraph*{The hidden abstraction.}
A close reading of this literature reveals an abstraction that is so widespread
that it rarely surfaces explicitly: measurement effort is operationalized as
\emph{the number of measured configurations}. Sarkar
\etal~\cite{sarkar2015} introduce a composite cost combining training
measurement cost, testing measurement cost, model-building cost, and prediction
error, then collapse acquisition effort to $2n$, treating each measured
configuration as having unit cost under a balanced train/test
split. Guo \etal~\cite{guo2013ASE} train predictors from samples of size $N$,
$2N$, $3N$, $M$ and discuss feasibility in terms of sample
cardinality. Westermann \etal~\cite{westermann2012} formulate efficiency as the
\emph{required number of measurement points} for a target error. Siegmund
\etal~\cite{siegmund2015FSE} report Pareto trade-offs between accuracy and
\emph{number of measurements}. Kaltenecker
\etal~\cite{kaltenecker2019} compare six sampling strategies at common sample
sizes. Pereira \etal{}'s systematic review~\cite{julianasurvey} confirms the
pattern: when ``cost'' is mentioned at all, it usually means cardinality of
measured samples.

This is the size-as-cost abstraction. It is convenient---sample size is
deterministic, comparable across systems, and easy to budget---but it bakes
in a strong assumption: that two sampling strategies selecting the same
number of configurations incur the same acquisition cost, or, in other
words, that configurations are equally expensive to measure. In practice,
configuration options can change per-configuration cost by orders of
magnitude. In x264, enabling slow analysis options or higher \texttt{ref}
frame counts dramatically inflates encoding time; in LLVM, certain
optimization passes turn a one-second compilation into a many-second one;
in archivers, compression level interacts non-linearly with input data and
with parallelism options. The same point holds for cost dimensions other
than time: when storage is the operational constraint, the cumulative
output size of measured configurations varies just as dramatically across
options. Strategies that preferentially explore some regions of the option
space (\eg solver-driven, distance-based, or coverage-based strategies)
may therefore systematically select cheaper or costlier configurations
than uniform random.

\paragraph*{A running example.}
A small synthetic configurable system makes the issue concrete.
\Cref{tab:running-example} shows a system with three binary options
$O_1, O_2, O_3$ and the eight configurations they induce. Each row
reports the elapsed time it takes to measure that configuration on a
fixed input. Concretely, $O_3$ is an option that activates expensive
analysis (think x264 high-quality motion estimation, or LLVM aggressive
optimization passes), so configurations with $O_3=1$ cost an order of
magnitude more to measure than configurations with $O_3=0$.

\begin{table}[t]
\centering
\caption{A synthetic configurable system with eight valid configurations,
three binary options, and per-configuration acquisition cost (here
instantiated as wall-clock measurement time, but the same logic applies
to any quantitative resource). Deliberately at toy scale to fit a single
page---the empirical study (\Cref{sec:experiment}) covers systems with
up to $\sim\!\!200{,}000$ valid configurations, $54$ options, and
training fractions from $1\%$ to $30\%$ of $|C|$, where the same effects
are amplified, not dampened. Two sampling strategies select the same
number of configurations ($n=3$), yet their empirical training costs
differ by an order of magnitude. After sampling, a regressor (random
forest, gradient-boosted trees, stepwise linear regression, etc.) is
trained on the selected (configuration, label) pairs and evaluated on
the remaining $5$ configurations; the resulting MRE is shown in the
bottom row.}
\label{tab:running-example}
\footnotesize
\setlength{\tabcolsep}{6pt}
\begin{tabular}{c c c c r c c}
\toprule
\multirow{2}{*}{Conf.} & \multicolumn{3}{c}{Options} & \multirow{2}{*}{$t$ (s)} & Strategy & Strategy \\
                       & $O_1$ & $O_2$ & $O_3$ &        & A   & B   \\
\midrule
$c_1$ & 0 & 0 & 0 & 1.0  & \checkmark &            \\
$c_2$ & 1 & 0 & 0 & 1.8  & \checkmark &            \\
$c_3$ & 0 & 1 & 0 & 2.2  & \checkmark &            \\
$c_4$ & 1 & 1 & 0 & 4.1  &            &            \\
$c_5$ & 0 & 0 & 1 & 6.5  &            & \checkmark \\
$c_6$ & 1 & 0 & 1 & 12.3 &            & \checkmark \\
$c_7$ & 0 & 1 & 1 & 14.0 &            & \checkmark \\
$c_8$ & 1 & 1 & 1 & 27.5 &            &            \\
\midrule
\multicolumn{4}{r}{$\SizeCost$ ($n$)}                       &      & \textbf{3}   & \textbf{3}    \\
\multicolumn{4}{r}{$\TrainCost$ (sum of $t$, in s)}         &      & \textbf{5.0} & \textbf{32.8} \\
\multicolumn{4}{r}{Predictor MRE on remaining $5$ configs}  &      & \textbf{20\%} & \textbf{10\%} \\
\bottomrule
\end{tabular}
\end{table}

\noindent
Suppose a practitioner has a budget of $n=3$ training configurations.
Strategy~A draws cheap configurations ($\{c_1,c_2,c_3\}$, all with
$O_3=0$); Strategy~B draws costly ones ($\{c_5,c_6,c_7\}$, all with
$O_3=1$). The conventional cost abstraction cannot distinguish them:
$\SizeCost=3$ either way. Empirically, Strategy~B costs more than six
times as much as Strategy~A.

A complete sampling+learning experiment does not stop at acquisition: it
trains a regressor on the $3$ measured configurations and evaluates its
MRE on the remaining five. Strategy~B's training set spans the expensive
$O_3=1$ regime that drives most of the cost mass, so its predictions
generalize better; Strategy~A's training set misses that regime
entirely. Bottom row of \Cref{tab:running-example}: an illustrative MRE
of $20\%$ for A and $10\%$ for B. Now the picture has two axes
simultaneously, and the cost-axis choice matters in a way that count
alone hides:
\begin{itemize}
    \item In the $(\SizeCost, \MRE)$ plane both strategies sit at $n=3$,
    with B at $10\%$ and A at $20\%$. \textbf{B dominates A under
    count-based cost}: same effort, lower error.
    \item In the $(\TrainCost, \MRE)$ plane A sits at $(5.0\,\mathrm{s},\,
    20\%)$ and B at $(32.8\,\mathrm{s},\, 10\%)$. \textbf{Neither
    dominates the other under empirical cost}: A pays $6.6\times$ less
    time to give up $10$ percentage points of MRE. A practitioner with a
    tight time budget might rationally pick A; one with a tight error
    budget might rationally pick B.
\end{itemize}
The Pareto status of A flipped---dominated under the size-as-cost
abstraction, Pareto-efficient under empirical cost---and the strategy you
would recommend depends on how heavily you weight time against accuracy
(RQ3) as well as on which $(s,n)$ pairs are on the frontier (RQ4).
Crucially, none of these conclusions could be drawn from sample size
alone. \Cref{app:running-example-figure} sketches the same comparison
visually, with sample sizes spanning $\{8,\ldots,128\}$ and Pareto fronts
under both cost models.

\Cref{tab:running-example} is deliberately small: three options, eight
configurations, $n=3$. That scale was chosen for readability, not for
realism. The phenomenon scales \emph{up}, not down. Real configurable
systems have between $11$ and $54$ binary options and between hundreds
and hundreds of thousands of valid configurations---LLVM at $1{,}024$,
LRZIP at $432$, x264 at $1{,}152$, BDB-C at $2{,}560$, HIPAcc at
$13{,}485$, 7Z at $68{,}640$, JavaGC at $\sim\!\!200{,}000$. With more
options, the cost-heterogeneity surface has more dimensions; with larger
configuration spaces, sampling strategies have more room to drift into
cost-biased subregions; and at the realistic training fractions used in
the literature (typically $1\%$ to $30\%$ of $|C|$, often $32$--$256$
configurations), small per-strategy biases compound across many
selections. Single options in real systems are known to multiply
per-configuration cost by orders of magnitude.

\begin{table}[t]
\centering
\caption{Order-of-magnitude scale check at x264 size: $|C|=1{,}152$
configurations over $16$ binary options, training fraction $n/|C|\approx
11\%$. The numbers are illustrative orders of magnitude derived from
plausible per-configuration encoding-time spreads, not measurements
(\Cref{sec:results} reports actual measured values). The qualitative
pattern of \Cref{tab:running-example} survives at scale: $\SizeCost$
tied across strategies while $\TrainCost$ differs by an order of
magnitude, and the Pareto status of the cheaper strategy still depends
on the cost dimension.}
\label{tab:running-example-scale}
\footnotesize
\setlength{\tabcolsep}{6pt}
\begin{tabular}{l r r}
\toprule
                                          & Strategy A    & Strategy B    \\
\midrule
$|C|$                                     & $1{,}152$     & $1{,}152$     \\
$n$ (training size)                       & $128$         & $128$         \\
sample fraction $\mathit{SizePct}=n/|C|$  & $\approx\!\!11\%$ & $\approx\!\!11\%$ \\
\midrule
$\TrainCost$                              & $\approx\!\!4$\,min  & $\approx\!\!40$\,min \\
$\mathit{TrainCostPct}=\TrainCost/\FullCost$ & $\approx\!\!3\%$  & $\approx\!\!33\%$ \\
calibration ratio $\Ratio$                & $\approx\!\!0.27$    & $\approx\!\!3.0$  \\
\midrule
predictor MRE                             & $\approx\!\!8\%$     & $\approx\!\!5\%$  \\
\bottomrule
\end{tabular}
\end{table}

\Cref{tab:running-example-scale} sketches the same comparison at x264
size. Replacing $|C|=8$ with $|C|=1{,}152$ changes the absolute numbers
but not the shape: a strategy can spend $\approx\!\!11\%$ of the
configuration budget while consuming $\approx\!\!33\%$ of the
measurement-time budget, a calibration ratio of $\approx\!\!3$ that the
size-as-cost abstraction makes invisible. The size and direction of that
gap is exactly what RQ1--RQ4 quantify in real systems.

The questions this paper asks (\Cref{sec:study_design}) are simply the
empirical version of what \Cref{tab:running-example} hints at: how
heterogeneous are real per-configuration costs (RQ1), does $n$ predict
$\TrainCost$ (RQ2), do strategy rankings change when $\TrainCost$
replaces $\SizeCost$ in the accuracy--cost trade-off (RQ3), do
Pareto-efficient $(s,n)$ pairs change under empirical cost (RQ4), and
does any of this depend on the system, the input, and the property being
predicted (RQ5)?

\paragraph*{What this paper does.}
We \emph{do not} claim that cost-aware sampling is novel. Sarkar
\etal~\cite{sarkar2015} already studied cost-efficient sampling. We claim
something narrower and, we believe, more defensible: \emph{the size-as-cost
abstraction itself has not been empirically validated}. We do not know, system
by system, whether $n$ is correlated with empirical measurement cost. We do not
know whether it is calibrated, \ie whether reducing $n$ by $x\%$ reduces
acquisition time by roughly $x\%$. We do not know whether different strategies
incur substantially different costs at fixed $n$. And we do not know whether
the cost-aware conclusions of prior work change when the cost term is replaced
by an empirical one.

The central question of this paper is therefore:
\begin{quote}
\emph{When is sample size a valid proxy for empirical measurement cost, when is
it miscalibrated, and does replacing count-based cost with empirical
measurement cost change the scientific conclusions previously drawn about
sampling strategies?}
\end{quote}
This framing is robust by design: if empirical cost closely tracks $n$ for the
systems and conditions we study, the abstraction is empirically justified
in those contexts and previous conclusions transfer; if it does not, the
abstraction hides strategy- and context-dependent effects that have to be
revisited.
This is not a blanket criticism of prior work. In configurable-systems
research, using sample size or number of measurements as an effort axis is a
practical and widespread convention in sampling, tuning, and transfer studies
\cite{kaltenecker2019,pereira2020,chen2021multiobj,chen2023aspirations,gong2024sempl}.
Some papers already take the stronger step of calibrating count budgets from
measured wall-clock time \cite{chen2021multiobj,chen2023aspirations}. Our claim
is narrower: to the best of our knowledge, prior work has not directly tested
whether sample size is itself a calibrated proxy for measured acquisition cost,
nor whether replacing count-based cost with measured cost changes rankings or
Pareto-efficient choices. We therefore position this paper as a validity study
of a common methodological assumption, not as a rejection of prior predictive
or optimization contributions.

\paragraph*{Approach.}
For each sampling strategy $s$, run $r$, training size $n$, system, input and
performance property, we compute an empirical \emph{training} measurement
cost
\[
\TrainCost(r,s,n)\;=\;\sum_{c\in S_{r,s,n}} \mathit{cost}(c)
\]
where $S_{r,s,n}$ is the training sample selected by the strategy. Crucially,
we restrict the cost term to training: validation/test measurements belong to
the experimental protocol, not to the strategy. We then compare
$\TrainCost$ to the count-based baseline $\SizeCost(s,n)=n$ along five axes:
heterogeneity at fixed $n$, correlation and calibration, ranking change under
combined accuracy--cost objectives, Pareto efficiency in the MRE--cost plane,
and stability across context. Datasets reuse benchmarks already used in prior
sampling studies---ten systems with measured per-configuration times,
plus the x264 case study with seventeen input videos and two performance
properties (encoding time and output size).

\paragraph*{Why this matters.}
Two practical consequences follow. First, if size is a poor proxy in even a few
of the systems and conditions we study, then the cost-aware claims that
underpin recommendations for practitioners (``strategy $X$ is cheaper than
strategy $Y$'', ``$n=2t$ is sufficient'') need to be re-stated in terms of
empirical time, and the conditions of validity need to be made explicit. Second,
if size is a good proxy under reasonably broad conditions, that is itself a
useful empirical finding: it tells the community that an abstraction it has
been using by default is not just convenient but actually justified for those
conditions. Either way, the size-as-cost question deserves to be answered with
data rather than assumed.

\paragraph*{Contributions.}
\begin{itemize}
    \item We define an empirical \emph{training} measurement cost,
    $\TrainCost(r,s,n)$, that captures the acquisition effort actually induced
    by a sampling strategy on a given system, input and property, and we
    distinguish it cleanly from validation/test cost so that the cost term
    reflects what a strategy controls.
    \item We empirically test, on ten configurable systems, whether sample size
    $n$ is (i) correlated with $\TrainCost$, (ii) calibrated as a budget proxy
    via the ratio $(\TrainCost/\FullCost)/(n/|C|)$, and (iii) free of
    strategy-specific bias.
    \item We compare strategy rankings under count-based and empirical-cost
    composite objectives across a sweep of trade-off parameters $\lambda$, and
    compute Pareto frontiers in the MRE--cost plane under both cost models.
    \item Using x264 with seventeen input videos and two performance properties
    (encoding time, output size) we characterize when the validity of the
    size-as-cost abstraction depends on context, including the asymmetric case
    in which the predicted property differs from the property that constitutes
    measurement cost.
    \item We pre-state twelve hypotheses (H1a--H5c) and report formal
    verdicts under explicit tests and thresholds (Kruskal--Wallis with tie
    correction and $\eta^2$ effect size, Spearman with Student-$t$
    $p$-values, percentile bootstrap CIs of the calibration ratio,
    Vargha--Delaney $\hat A_{12}$), giving a per-system and per-context
    picture of which claims are supported, refuted, or inconclusive
    (\Cref{tab:hypothesis-verdicts}).
    \item We release the experimental pipeline, raw measurements, and analysis
    artifacts at \red{anonymized URL}.
\end{itemize}

\paragraph*{Roadmap.}
\Cref{sec:background} discusses configurable systems, sampling strategies, and
prior cost models. \Cref{sec:cost_model} formalizes the empirical training
cost, the count-based baseline, the composite objectives, and the Pareto
formulation. \Cref{sec:study_design} states the five research questions, the
hypotheses, and the statistical analysis. \Cref{sec:experiment} details
datasets, strategies, learners, sample sizes, and evaluation
protocol. \Cref{sec:results}--\Cref{sec:discussion} report and discuss the
results. \Cref{sec:threats} addresses threats to
validity. \Cref{sec:related_work} positions our work in the broader literature,
and \Cref{sec:conclusion} concludes.
